\section{From Conventional Reactor Control to Nuclear-State Control}
\label{sec:control_framework}

Conventional reactor control is based on the manipulation of macroscopic variables that influence neutron balance. At the simplest level, the evolution of a multiplying system may be represented through the effective multiplication factor $k_{\mathrm{eff}}$, or equivalently through the reactivity
\begin{equation}
\rho
=
\frac{k_{\mathrm{eff}}-1}{k_{\mathrm{eff}}}.
\end{equation}

For a fissile population $N_f$, the neutron-induced fission rate can be written schematically as
\begin{equation}
R_f(t)
=
N_f
\int
\sigma_f(E)\,
\Phi(E,t)\,
dE.
\label{eq:fission_rate_control}
\end{equation}

In this formulation, the main time-dependent quantity accessible to conventional control is the neutron field. Control rods, soluble absorbers, moderator conditions, coolant density, fuel temperature, and geometry ultimately modify neutron production, absorption, moderation, or leakage.

A simplified conventional control chain may therefore be represented as
\begin{equation}
u_C(t)
\longrightarrow
\mathcal{A}_C
\longrightarrow
\Phi(E,t)
\longrightarrow
R_f(t),
\end{equation}
where $u_C(t)$ denotes a conventional command and $\mathcal{A}_C$ the associated physical actuator.

The nuclear properties entering $\sigma_f(E)$ are not normally treated as dynamically controllable quantities. For a given isotope and thermodynamic regime, they are effectively constitutive properties encoded in evaluated reaction data.

The present work investigates the consequences of relaxing this assumption.

\subsection{Nuclear state as an explicit variable}

Let the state of the fissile nuclear ensemble be represented by a density operator
\begin{equation}
\rho_N.
\end{equation}

A density operator is preferable to a pure state because macroscopic fissile material will generally contain an ensemble of nuclei occupying multiple states and interacting with a complex environment.

The conventional fission cross section may then be regarded as an ensemble-averaged response associated with an uncontrolled or reference distribution,
\begin{equation}
\sigma_f(E)
=
\sigma_f(E;\rho_N^{(0)}).
\end{equation}

If an external interaction changes the state distribution,
\begin{equation}
\rho_N^{(0)}
\rightarrow
\rho_N(t),
\end{equation}
the more general response is
\begin{equation}
\sigma_f
=
\sigma_f(E;\rho_N).
\end{equation}

The fission rate becomes
\begin{equation}
R_f(t)
=
N_f
\int
\sigma_f(E;\rho_N(t))
\Phi(E,t)\,dE.
\label{eq:fission_rate_state}
\end{equation}

Equation~\eqref{eq:fission_rate_state} contains two conceptually distinct control channels:
\begin{equation}
\Phi(E,t)
\rightarrow
R_f(t)
\end{equation}
and
\begin{equation}
\rho_N(t)
\rightarrow
\sigma_f^{\mathrm{eff}}(t)
\rightarrow
R_f(t).
\end{equation}

Conventional control modifies the neutron environment. Nuclear-state control would instead modify the response of the fissile material itself.

\subsection{A second control channel}

We introduce a generic external control variable $u_Q(t)$ representing an interaction capable of modifying the nuclear state:
\begin{equation}
u_Q(t)
\longrightarrow
\rho_N(t)
\longrightarrow
\sigma_f^{\mathrm{eff}}(t)
\longrightarrow
R_f(t).
\label{eq:q_control_chain}
\end{equation}

Its physical implementation is deliberately left unspecified. It could correspond to a coherent electromagnetic field, an electron-mediated process, an X-ray or gamma-ray pulse sequence, or another externally tunable interaction.

The complete controlled fission system may therefore possess two parallel pathways,
\begin{align}
u_C(t) &\longrightarrow \Phi(E,t),\\
u_Q(t) &\longrightarrow \rho_N(t),
\end{align}
which both influence $R_f(t)$.

\subsection{State-dependent fission probability}

Consider a set of nuclear states $\{|i\rangle\}$, each characterized by a state-dependent fission response $\sigma_f^{(i)}(E)$. For an incoherent mixture with populations $p_i$,
\begin{equation}
\sigma_f^{\mathrm{eff}}(E)
=
\sum_i
p_i
\sigma_f^{(i)}(E).
\label{eq:incoherent_sigma}
\end{equation}

Thus, external redistribution of state populations,
\begin{equation}
p_i^{(0)}
\rightarrow
p_i(t),
\end{equation}
is sufficient to modify the effective fission response if at least two states have unequal fission cross sections.

A more general quantum description contains off-diagonal density-matrix elements $\rho_{ij}$ with $i\neq j$. Formally, the fission probability can be written as
\begin{equation}
P_f(E)
=
\mathrm{Tr}
\left[
\rho_N\,
\hat{\Pi}_f(E)
\right],
\label{eq:fission_operator}
\end{equation}
where $\hat{\Pi}_f(E)$ denotes an effective operator associated with evolution into the fission channel.

Equation~\eqref{eq:fission_operator} expresses the central idea compactly: if an external control changes $\rho_N$ and the fission operator is state sensitive, then the fission probability is externally modifiable.

\subsection{From microscopic control to macroscopic reactivity}

Let
\begin{equation}
\delta\sigma_f
=
\sigma_f^{\mathrm{controlled}}
-
\sigma_f^{(0)}
\end{equation}
be the state-induced modification of the microscopic fission response. At the macroscopic level,
\begin{equation}
\delta\Sigma_f
=
N_f\,\delta\sigma_f.
\end{equation}

A nonzero $\delta\Sigma_f$ perturbs neutron production and therefore, in principle, the multiplication factor:
\begin{equation}
\delta\rho_N
\rightarrow
\delta\sigma_f
\rightarrow
\delta\Sigma_f
\rightarrow
\delta k_{\mathrm{eff}}
\rightarrow
\delta\rho.
\end{equation}

For bookkeeping, we introduce a hypothetical effective reactivity contribution $\rho_Q$ and write, at first order,
\begin{equation}
\rho_{\mathrm{tot}}
=
\rho_{\mathrm{conv}}
+
\rho_Q.
\label{eq:rho_decomposition}
\end{equation}

Equation~\eqref{eq:rho_decomposition} is not assumed to hold as an exact additive decomposition in arbitrary nonlinear regimes; it is a perturbative definition that permits direct comparison with conventional reactivity effects.

\subsection{Necessary physical conditions}

For nuclear-state control to constitute a meaningful additional degree of freedom, several conditions must hold:
\begin{enumerate}
\item an external interaction must reproducibly modify the nuclear-state distribution,
\begin{equation}
u_Q(t)\rightarrow\delta\rho_N(t);
\end{equation}
\item the relevant states must have distinguishable fission responses,
\begin{equation}
\frac{\partial P_f}{\partial\rho_N}\neq0;
\end{equation}
\item the state modification must survive long enough to affect the relevant fission interaction;
\item a sufficiently large fraction of nuclei must be addressed;
\item the ensemble-averaged change must survive competing relaxation, statistical fluctuations, and ordinary reactor feedback.
\end{enumerate}

The proposed mechanism can therefore be summarized by
\begin{equation}
\boxed{
u_Q
\rightarrow
\delta\rho_N
\rightarrow
\delta P_f
\rightarrow
\delta\Sigma_f
\rightarrow
\delta k_{\mathrm{eff}}
}.
\label{eq:necessary_chain}
\end{equation}

Every arrow in Eq.~\eqref{eq:necessary_chain} is a separate physical problem. The following sections examine which links are already supported by existing nuclear physics and which remain open.
